\section{Discussion}

\subsection{What the evidence supports}

The compressibility result is robust. Across two backbone families, five
seeds, and a 32-fold range of embedding dimension, macro AUC varies by less
than the seed-to-seed standard deviation of the training procedure itself, and
an equivalence test with a pre-registered margin confirms this rather than
merely failing to reject a null. For twelve-lead resting ECG and five-class
superdiagnostic labels, the diagnostic manifold is low-dimensional enough that
sixteen coordinates suffice.

The practical consequence is genuine but narrower than the compression
literature typically advertises. A single artefact replaces a family of
per-dimension models, with the corresponding reduction in training,
validation, storage and---in a regulated setting~\cite{fda2021}---certification
burden. Where embeddings are transmitted or indexed at scale, the 32-fold
storage reduction is the dominant saving.

\subsection{What it does not support}

\textbf{Inference compute.} Nesting leaves it untouched, as
Section~\ref{sec:compute} measures directly. This is inherent to the method
rather than an implementation shortcoming: prefix truncation happens after the
encoder has run. Papers that present MRL as an edge-efficiency technique
without this qualification are, we think, misleading readers about where the
saving lies. The complementary axis---quantisation, distillation, a smaller
encoder---remains entirely available and is where compute savings must come
from. We measure the first of these directly (Section~\ref{sec:quant}): INT8
quantisation composes with nesting, and the differential penalty column of
Table~\ref{tab:quant} reports whether it degrades small prefixes more than
large ones. That interaction has not to our knowledge been examined for nested
representations in any domain, and it is the one that determines whether the
two techniques can be stacked safely at the smallest operating points.

\textbf{Wearable deployment.} PTB-XL is clinical, twelve-lead, resting and
wet-electrode. Our reduced-lead and artefact experiments probe the gap but do
not close it. Establishing wearable viability requires ECGs acquired by
wearable hardware, ideally with paired clinical reference, and execution on the
target microcontroller. Neither is present here, and a device-tier table
mapping dimensions to smartwatches would misrepresent what has been shown. We
have removed it.

\textbf{Diagnostic accuracy.} Our absolute performance sits below the
strongest published PTB-XL systems. The mechanism studied here is orthogonal
to backbone quality, and the flat accuracy--dimension profile is measured
within a fixed training protocol, but a paper proposing nesting for clinical
use would need to demonstrate it on a competitive encoder.

\subsection{On the interpretation of flatness}

It is worth being precise about why performance is flat, because two
explanations are easily conflated and they have different implications.

Under the \emph{regularisation} account, the multi-granularity loss constrains
the encoder in a way that improves the leading coordinates. Under the
\emph{low-dimensionality} account, the five-superclass task simply does not
require more than a few coordinates, and any reasonable training procedure
concentrated on a prefix would find them. The two are distinguishable: the
regularisation account predicts that MRL beats same-backbone fixed-dimension
training at small $d$, whereas the low-dimensionality account predicts
equivalence. The matched baselines of Table~\ref{tab:main} are precisely this
test, and the probing analysis of Section~\ref{sec:probing} bears on it
further---if all descriptors saturate immediately, low-dimensionality is the
parsimonious reading.

The previous version asserted the regularisation account on the basis of a
cross-architecture comparison, which cannot distinguish the two. We flag this
explicitly because the same inferential slip is common in the adaptive-model
literature: an adaptive model outperforming a differently-architected fixed
model is evidence about architectures, not about adaptivity.

The observation that $d$=16 sometimes slightly \emph{exceeds} $d$=512 warrants
similar restraint. The gap is well inside seed noise, and treating it as a
real regularisation effect over-reads the data. Our tables report it with
uncertainty attached and draw no mechanism from it.

\subsection{Limitations}

\begin{enumerate}
\item \textbf{Task granularity.} Only the five-class superdiagnostic task is
evaluated. The 44-class diagnostic and 71-class ``all'' PTB-XL tasks have
higher intrinsic dimensionality, and there is no reason to expect $d$=16 to
suffice for them. If the low-dimensionality account is correct, the
sufficient prefix should grow with task granularity---a prediction the
framework makes and that we have not tested.

\item \textbf{No wearable hardware.} No microcontroller execution, no
wearable-acquired data. Latency on embedded targets is unmeasured and, given
the analysis above, unlikely to benefit from nesting in any case.

\item \textbf{Label mapping in transfer.} SNOMED-CT to superclass mapping is
many-to-one and lossy; external results measure concept transfer under shift,
not identical-task performance.

\item \textbf{Synthetic artefacts.} Injected noise at controlled SNR
approximates but does not reproduce dry-electrode ambulatory recording.

\item \textbf{Scale.} PTB-XL is small by contemporary standards. Corpora such
as MIMIC-IV-ECG and CODE-15\% offer two to three orders of magnitude more
data; whether nesting behaves the same at that scale, and whether the
sufficient prefix grows with dataset size, is open. Our pipeline supports
these corpora but they require credentialed access we have not exercised.

\item \textbf{Single sampling rate.} Only the 100\,Hz variant is evaluated;
higher rates resolve morphology that may not compress identically.
\end{enumerate}

\subsection{Future work}

The most valuable next step is the one we could not take: wearable-acquired
ECGs with clinical reference labels, evaluated on target hardware. Beyond
that, three directions follow from the analysis. Combining nesting with
quantisation would address the compute axis that nesting leaves untouched.
Extending to finer diagnostic granularities would test whether the sufficient
prefix scales with task complexity, which is the sharpest prediction the
low-dimensionality account makes. Applying the same probing methodology to
other biosignal modalities (EEG, PPG, EMG) would establish whether nested
compressibility is a property of ECG specifically or of physiological
time series generally.
